% Generated draft: requires mathematical and source review.
\documentclass[11pt]{article}
\usepackage[margin=25mm]{geometry}
\usepackage{fontspec}
\setmainfont{Latin Modern Roman}
\usepackage{amsmath,amssymb,mathtools,unicode-math}
\setmathfont{Latin Modern Math}
\usepackage{enumitem,needspace,xurl,hyperref}
\setlength{\parindent}{0pt}
\setlength{\parskip}{5pt}
\setlength{\emergencystretch}{3em}
\newcommand{\N}{\mathbb{N}}
\newcommand{\Z}{\mathbb{Z}}
\newcommand{\Q}{\mathbb{Q}}
\newcommand{\R}{\mathbb{R}}
\newcommand{\C}{\mathbb{C}}
\newcommand{\F}{\mathbb{F}}
\newcommand{\A}{\mathbb{A}}
\newcommand{\PP}{\mathbb{P}}
\newcommand{\E}{\mathbb{E}}
\newcommand{\eps}{\varepsilon}
\newcommand{\dd}{\,\mathrm{d}}
\begin{document}
\section*{UnsolvedMath ID: 11000117}
\textbf{Source problem code:} catalogue entry
\subsection{Definitions and mathematical statement}
For a positive integer \(r\), let
\[
 \Gamma_r=\langle T_c^r:c\text{ is a simple closed curve defining a Dehn twist in }\operatorname{Mod}_g\rangle.
\]
This is the source's normal subgroup generated by the \(r\)-th powers of all Dehn twists, including all curve types intended by that surface convention. For a fixed nonseparating curve \(a\), set
\[
 N_r=\langle\!\langle T_a^r\rangle\!\rangle_{\operatorname{Mod}_g}.
\]
The genus range and surface convention are not stated in the excerpt.

\emph{Problem 2.12.} Is \([\operatorname{Mod}_g:\Gamma_r]=\infty\)? Determine the answer for the intended genera and positive exponents, retaining the known exceptions reported in the source. It attributes to Humphries [17] that \(N_r\) has finite index for \(r=2\) in every intended genus and for \((g,r)=(2,3)\). Since \(N_r\subseteq\Gamma_r\), these give negative answers to the infinite-index question in those cases.

It further states, citing [17] and [10], that \(N_r\) has infinite index if \(r\notin\{2,3,4,6,8,12\}\). This statement requires an omitted range qualification: as written it includes \(r=1\), whereas in the usual mapping-class setting nonseparating twists normally generate the group. The draft preserves the reported exclusion set but does not treat that unrestricted claim as verified. Also, infinite index of the smaller subgroup \(N_r\) alone does not imply infinite index of \(\Gamma_r\); the two subgroups must not be identified.

The extraction then gives separate context: \(\operatorname{PMod}_{g,p}^{b}\) consists of mapping classes fixing each puncture, and for \(g\geq4\) it states
\[
 H_2(\operatorname{PMod}_{g,p}^{b};\Z)\cong\Z^{p+1},
\]
citing [13], [39], [33], with the warning that [33] interchanges the subscript \(p\) and superscript \(b\). The exact boundary/puncture conventions and the scope of that homology statement must be checked. This is adjacent context, not a further condition on the power-subgroup question.
\subsection{Short English statement}
Ask when the subgroup generated by all \(r\)-th powers of Dehn twists has infinite index. The source gives finite-index exceptions via the normal closure of one nonseparating twist, but its other normal-closure range has a missing \(r=1\) qualification and does not itself settle the larger subgroup.
\subsection{Sources}
\begin{itemize}
\item Local UnsolvedMath registry, permanent ID 11000117, source code catalogue entry, Problem 2.12, PDF page 97. No individual proposer is supplied. Page distinguishes this from other Problem 2.12 headings in the volume.
\item Farb, \emph{Problems on Mapping Class Groups and Related Topics}. Supplied URL: \url{https://www.math.uchicago.edu/~farb/papers/mcgbook.pdf}; recorded access 2026-07-29; extraction pdf-font-structure. The edited volume contains 329 explicitly typeset headings across contributed chapters.
\item Humphries [17] and reference [10] are cited for the normal-closure index statements. References [13], [39], [33] are cited for the adjacent pure-mapping-class homology claim; the source warns of interchanged indices in [33]. Full bibliographic entries are not supplied.
\item Supplied status open and OPEN-TRIAGE. The summary says only ``Truncated.'' Current primary-literature status and redistribution terms for the publicly accessible author/publisher copy remain NEEDS\_REVIEW.
\end{itemize}
\end{document}
